% ======================= 译 者 序 =======================
\clearpage
\chapter*{译\ 者\ 序}
\markboth{译者序}{译者序}
\addcontentsline{toc}{chapter}{译者序}

\noindent
二〇二一年十二月，笔者赴莫斯科大学历史系进行为期三周的学术访问。访问期间，巴兰尼科夫教授在其办公室出示了一只铁盒，盒盖上的火漆已经开裂，编号清晰可辨，为 \textsc{ф.\,117, оп.\,4, д.\,2187}。据其介绍，盒内材料在国内尚无公开著录。\footnote{苏联科学院远东委员会（1949—1991）的档案于一九九二年部分解密。第 117 全宗第 4 目第 2187 号卷宗至今未见于任何公开目录，其存在仅由两份互相引用的移交清单间接证明。参见俄罗斯科学院档案馆编《解密档案目录补编（1992—2010）》，莫斯科，2011 年，第 288 页。}

盒内计有打字稿三十七份、缩微胶片四卷，以及金属圆片一枚，直径不足两厘米，表面刻有极细的同心圆纹，肉眼几乎不可见。笔者最初的判断是这类材料多属伪托，不足采信。三个月后，莫斯科大学地质系对该圆片作了铼—锇同位素稀释质谱分析，其 \textsuperscript{187}\textsc{re}/\textsuperscript{187}\textsc{os} 比值为 3.147±0.006，与阿波罗十七号所采月海高地玄武岩的测定值在误差范围之内完全一致。\footnote{分析报告编号 МГУ-ГЕО-2022-014，结论部分有一句被作者划去的话：“建议不再重复采样。”}

这一结果使笔者不得不重新考虑原有的判断。本书所呈现的，即是对这批材料及其相关问题的一次系统整理。

\section*{一、\ 研究方法上的差异}

巴兰尼科夫教授的治学路径，与国内中国史研究的通行做法有明显不同。他并不以考据为主，尽管对《旧唐书》《明实录》等书的熟习程度不亚于专门研究者；他也不以辨伪为务，尽管其一生工作的大部分都涉及材料的真伪判断。他所做的工作，他称之为“重构”，即从一个文明留下的全部物质与文本痕迹出发，反推其曾经拥有、后来又被放弃的技术条件。

用他自己的话说，历史学家不应只问“这不可能”，而应当问“要让它可能，需要多少能量”。\footnote{巴兰尼科夫《重构史学导论》，莫斯科大学出版社，2015 年，第 9 页。该页页边另有手书一句：“此句被引用的次数，远超我被阅读的次数。”}

这一方法被他概括为“四重证据法”，即文献、器物、同位素与轨道四项。前三项国内学者并不陌生，王国维的“二重证据法”已开其端；第四项则为其独有。他认为，任何大规模的技术活动都会在近地轨道留下不易磨灭的残骸层，而这一层“轨道考古学”的记录，比任何正史都更难篡改。\footnote{据欧洲空间局 2019 年发布的轨道碎片统计，地球轨道上可追踪的十厘米以上人造物体约三万四千件，其中约百分之零点三来源不明。巴兰尼科夫在全书各处反复引用这一比例，并称之为“我们唯一的诚实证人”。}

正是借助这一方法，他从莫斯科档案馆未公开的冷战时期考古卷宗中，读出了国内学界未曾设想的内容，包括一个被反复重启的文明、一段被系统删除的技术史，以及一场在正史中仅占两行字的、星际规模的冲突。

\section*{二、\ 本书结论的评价问题}

本书的核心结论是惊世骇俗的，笔者不打算在序言中为其辩护。

一个中国历史研究者初读第一章，其反应大概会与笔者初见那枚圆片时相同，即先觉其荒谬，继而觉其唐突，最后归于沉默。所以然者，在于本书所使用的材料，其卷次、页码、版本乃至避讳缺笔均准确无误；被替换掉的只是材料与材料之间的空白。

这一点是本书最令人不安之处。它几乎全部建立在真实的史料之上，只是把史料之间的空白，用一种难以接受、却也难以证伪的方式填满了。\footnote{关于这种以真实材料支撑不可证结论的写法，可参见张建国《伪史学的技术伦理》，《史学理论研究》2022 年第 3 期，第 55—71 页。该文认为此类写作“在方法上无可指责，在结论上不可接受”，并建议“将其归入文学”。}

中国社会科学院历史研究所伪史学研究中心之所以决定承担此书的翻译，理由正在于此。一种谬误，若其每一个步骤都合乎规范，则它已不再是简单的谬误，而是一份关于规范本身的诊断书。翻译此书，并非因为相信其结论，而是为了弄清一种如此严整的荒谬究竟是如何被建造出来的。

\section*{三、\ 翻译体例说明}

其一，全书俄文原注保留原有格式，中文补注以“译者注”标明，不作合并。原注中大量使用的缩略语（如 \textsc{цга рф}、\textsc{ив ран}）按俄罗斯档案学界通例译出，首次出现时附原文。

其二，作者惯用的一组术语，中译取其字面义而非通行义。例如 \ru{перезапуск} 译作“重启”而不译作“复兴”，\ru{поле забвения} 译作“记忆消除场”而不译作“遗忘场”。这一处理曾与作者本人通信商定，他在回信中说：“随您怎么译。反正它们都会被删掉。”\footnote{作者 2022 年 3 月 14 日致译者电子邮件，现存译者处。此句之后尚有一行，因涉及未刊史料，遵作者嘱暂不发表。}

其三，本书第三版较前两版增补了第一章第三节的全部内容，并恢复了 1979 年被删去的、关于大气适应性退化的两页半文字。据作者称，这两页半是理解全书的枢纽。\footnote{删节本的影印件见巴兰尼科夫《被删去的两页半》，载《远东问题》1993 年第 2 期。该刊同期编者按称：“本刊发表此文，不代表本刊同意其结论。”}

其四，书中所有插图均为作者据档案原件重绘。原件照片因技术原因未能随书出版，此事已多次要求说明，未获答复。

\section*{四、\ 一点说明}

翻译此书的过程中，笔者曾反复考虑一个问题，即一部在方法上无可指摘、在结论上难以接受的著作，应当以何种方式呈现给中文读者。

笔者的结论是，不作任何调和。凡作者之判断，一律照译；凡笔者之保留，一律入注。读者所见，即为原貌。

至于这些判断究竟成立与否，须由读者自行裁断。本书所提供的，不是答案，而是一份完整的、可供检验的材料。

\vspace{10mm}
\begin{flushright}
{\small 李\ 文\ 亚}\\[2pt]
{\footnotesize\color{pencil} 二〇二三年九月，于北京建国门内大街五号}
\end{flushright}

% ======================= 中文版序言 =======================
\clearpage
\chapter*{中文版序言}
\markboth{中文版序言}{中文版序言}
\addcontentsline{toc}{chapter}{中文版序言}

\noindent
本书的中译者在序言中写道，我的结论“惊世骇俗”。这一说法客气得近乎虚伪。更准确的表述是：我的结论之所以令人不适，并非因为它们违背了史料，而是因为它们违背了诸位对史料的期待。

中国的史料是世界上最完备的史料，这一点我从不否认，并且由衷敬佩。但也正因如此，中国的史学传统形成了一种可以称为文本崇拜的习惯：凡不见于文字者，即视为不存在；凡见于文字者，即视为曾经发生。这是一种极为经济的世界观，它把整个宇宙压缩进一部可以随身携带的丛书之中。\footnote{作者此处原有一长注，列举中国类书对异域记载的系统性简化，后整段删除。据其助手称，删去的理由是“太刻薄了，而且他们译不过来”。}

各位即将读到的四章，分别处理四个时间断面。它们并非四个孤立的奇谈，而是同一条曲线的四个采样点，即一个文明如何获得技术，如何滥用技术，如何被封存，以及如何被自己的封存所出卖。

在此预先声明三点。

第一，本书不提供证据链，只提供证据点。点与点之间的连线是我的判断，凡判断皆可能有误。我欢迎反驳，但请反驳连线，不必反驳那些点；那些点是我用三十一年时间从十一个国家的档案中逐一辑出的。\footnote{作者此处用了一个俄语俗语 \ru{до мозолей}（直到磨出茧子），译者按字面译出。}

第二，本书不承认民族感情可以作为学术标准。若某一结论令某一民族感到不快，那只能说明历史本身曾经令人不快，而不能说明结论有误。历史学不是安慰剂。

第三，也是最重要的一点：本书第三版恢复了 1979 年被删去的两页半。那两页半讨论的是一个技术细节，即大气密度梯度对减速装置的影响系数。这一系数，以最朴素的方式证明了：我们的祖先不是从树上下来的，是从天上掉下来的。

这句话，我在莫斯科讲了四十三年。今天，它第一次被译成中文。

愿它被读完，而不只是被引用。

\vspace{10mm}
\begin{flushright}
{\small А.\,В.\,巴兰尼科夫}\\[2pt]
{\footnotesize\color{pencil} 二〇二三年十一月，于莫斯科列宁山}
\end{flushright}
